\documentclass{jarticle}
\usepackage[dvipdfmx]{graphicx} % Required for inserting images (platex + dvipdfmx でPNGを読むため)
\usepackage{amsmath} % For \tag command
  \graphicspath{{figs/}} % 図は figs/workN/ から読む（Overleaf には figs フォルダごとアップロード）
% 図のパネル: \panel{画像}{記号}  （1つの figure 環境に複数の図を (a)(b)... で並べる．第1章 1.2.5）
\newcommand{\panel}[2]{%
  \begin{minipage}[t]{0.48\linewidth}\centering
    \includegraphics[width=\linewidth,bb=0 0 800 600]{#1}\\ {\small (#2)}
  \end{minipage}}

\begin{document}

\begin{center}

    {\huge  数理工学実験 レポート} \\[2cm]
    {\LARGE べき乗法およびQR法を用いた固有値計算} \\[2cm]

    \begin{large}
    \begin{tabular}{ll}
        {実験実施年月日：} & 2026年10月06日 \\
        {実験実施場所：} &  工学部総合校舎数理計算機室\\
        {レポート提出年月日：} & 2026年10月26日 \\[1cm]
        学年・学生番号： & 工学部　2回生 \\
                       & 1029379268 \\
        氏名： & 脇田悠希
    \end{tabular}
    \end{large}

\end{center}

\newpage

\section{はじめに}
数値線形代数の主要テーマである連立 1 次方程
式や行列固有値問題は、科学技術計算から日常生
活まで、ありとあらゆるところに現れる基本問題
である。数値線形代数の基礎的な知識と技術を身につけることを目的として、 連立 1
次方程式
Ax = b (2.1)
および行列固有値問題
Ax = λx (2.2)
の数値解法を実装し，数値実験によりその性能を
調べる

\section{原理・方法}
\subsection{消去法}
Gaussの消去法では、
\begin{equation}
    A x = b
    \tag{2.1}
\end{equation}
の解を次のように求める。

(1) 拡大係数行列$(A,b)$を行基本変形により行階段形$(A\prime,b\prime)$にする。
(2) 方程式$A\prime x=b\prime$の解を下の漸化式\eqref{eq:backsub}により求める。
\begin{equation}
    x_i = \frac{1}{a'_{i,i}} \left( b'_i - \sum_{j=i+1}^{n} a'_{i,j} x_j \right)
    \label{eq:backsub}
\end{equation}
なお、行基本変形は以下。
\subsection{前進消去による行基本変形}
次の手続きを $k = 1, 2, . . . , n-1$ の順に繰り
返して拡大係数行列 $(A, b)$ を逐次更新する．

(1) 第 $k$ 列の下三角成分 $a_{k,k}, ..., a_{n,k}$ の中か
らノルムが最大のものをひとつ選ぶ．（この $a_{i,k}$ をピ
ボットという．）

(2) $i \neq k$ ならば第 $k$ 行と第 $i$ 行を入れ替える．

(3)（$a_{k+1,k}, ... , a_{n,k}$ がすべて $0$ になるように）各 $k + 1 \leq i \leq n$に対して第 $i$ 行から第 $k$行の$\frac{a_{i,k}}{a_{k,k},}$倍を引く．

\subsection{LU分解による解法}
以下の各手法の原理は文献\cite{text}による。
行列$A$を、単位下三角行列（対角成分がすべて1の下三角行列）$L$と上三角行列$U$の積$A=LU$に分解することをLU分解という。$A$が正則ならば、行の並べ替えを表す置換$\pi$とその置換行列$P_\pi$を用いて$P_\pi A=LU$と分解できる。これは前進消去で行う行基本変形を行列として記録したものであり、$\pi$を恒等置換として、$k=1,\ldots,n-1$の順に次を繰り返して求められる。

(1) 第$k$列の$a_{k,k},\ldots,a_{n,k}$のうち絶対値が最大の$a_{i,k}$をピボットとして選ぶ。

(2) $i\neq k$ならば$A$と$L$の第$k$行と第$i$行を入れ替え、$\pi$に左から互換$(k\ i)$を掛ける。

(3) $L$の第$k$列を$(0,\ldots,0,1,a_{k+1,k}/a_{k,k},\ldots,a_{n,k}/a_{k,k})^{\mathsf T}$とする。

(4) 各$k+1\leq i\leq n$について、$A$の第$i$行から第$k$行の$a_{i,k}/a_{k,k}$倍を引く。

最後の$A$が$U$である。$P_\pi A=LU$のとき、$Ax=b$は$Ly=P_\pi b$、$Ux=y$と同値であるから、
\begin{align*}
    y_i&=b_{\pi(i)}-\sum_{j=1}^{i-1}\ell_{i,j}y_j & (i&=1,\ldots,n),\\
    x_i&=\frac{1}{u_{i,i}}\Bigl(y_i-\sum_{j=i+1}^{n}u_{i,j}x_j\Bigr) & (i&=n,\ldots,1)
\end{align*}
により$y$、$x$を順に求める（前進代入と後退代入）。

\subsection{べき乗法}
$A$の固有値を$|\lambda_1|>|\lambda_2|>\cdots>|\lambda_n|$と並べ、対応する固有ベクトルを$v_1,\ldots,v_n$とする。ベクトル$x$を$x=\sum_{i=1}^{n}c_iv_i$と展開すると
\[
    A^kx=\lambda_1^k\Bigl(c_1v_1+\sum_{i=2}^{n}c_i\Bigl(\frac{\lambda_i}{\lambda_1}\Bigr)^kv_i\Bigr)
\]
であり、$|\lambda_i/\lambda_1|<1$（$i\geq2$）であるから、$c_1\neq0$ならば$A^kx$の向きは$v_1$に近づく。これを用いて第1固有値と第1固有ベクトルを求めるのがべき乗法である。初期ベクトル$x_0$から$k=0,1,2,\ldots$の順に次を繰り返す。

(1) $y_k=Ax_k$とする。

(2) $x_k$の成分のうち絶対値が最大のものの番号を$i$として、$\mu_k=y_{k;i}/x_{k;i}$とする。

(3) $x_{k+1}=y_k/\|y_k\|_2$とする。

(4) $|\mu_k-\mu_{k-1}|/|\mu_k|<\varepsilon$ならば終了し、$\mu_k$を$\tilde{\lambda}_1$、$x_{k+1}$を$\tilde{v}_1$とする。

収束は収束率$|\lambda_2|/|\lambda_1|$の1次収束である。

\subsection{逆反復法}
$\sigma$が$A$の固有値でないとき、$(A-\sigma I)^{-1}$の固有値は$1/(\lambda_i-\sigma)$であり、その絶対値が最大となるのは$\sigma$に最も近い$\lambda_i$に対するものである。したがって$(A-\sigma I)^{-1}$にべき乗法を適用すると、$\sigma$に最も近い固有値$\lambda_i$とその固有ベクトルが求まる。これが逆反復法である。各反復では$y_k=(A-\sigma I)^{-1}x_k$を計算する代わりに連立1次方程式$(A-\sigma I)y_k=x_k$を解く。係数行列は反復の間変わらないので、最初に1回だけLU分解しておけば、各反復は前進代入と後退代入だけで済む。べき乗法と同様に$\mu_k$を求め、$\sigma+1/\mu_k$を$\lambda_i$の近似値とする。収束率は$|\lambda_i-\sigma|/|\lambda_j-\sigma|$（$\lambda_j$は$\sigma$に2番目に近い固有値）であり、$\sigma$が$\lambda_i$に近いほど速く収束する。

\subsection{QR分解}
行列$A=(a_1,\ldots,a_n)$を、直交行列$Q$（$Q^{\mathsf T}Q=I$）と上三角行列$R$の積$A=QR$に分解することをQR分解という。本実験では修正Gram--Schmidt法を用いた。$j=1,\ldots,n$の順に、$b_{1,j}=a_j$から
\begin{align*}
    r_{i,j}&=(q_i,b_{i,j}),\quad b_{i+1,j}=b_{i,j}-r_{i,j}q_i \qquad (i=1,\ldots,j-1),\\
    r_{j,j}&=\|b_{j,j}\|_2,\quad q_j=\frac{b_{j,j}}{r_{j,j}}
\end{align*}
として、$Q=(q_1,\ldots,q_n)$と$R=(r_{i,j})$を求める。

\subsection{QR法}
$A_0=A$から$k=0,1,2,\ldots$の順に、$A_k=Q_kR_k$とQR分解し$A_{k+1}=R_kQ_k$とする反復をQR法という。$R_k=Q_k^{\mathsf T}A_k$より$A_{k+1}=Q_k^{\mathsf T}A_kQ_k$であるから、各$A_k$は$A$と同じ固有値をもつ。$A_k$は上三角行列に近づき、その対角成分には固有値が絶対値の大きい順に並ぶ。狭義下三角成分$a_{k;i,j}$（$i>j$）は収束率$|\lambda_i|/|\lambda_j|$で0に1次収束するので、全体の収束の速さは$\max_i|\lambda_{i+1}|/|\lambda_i|$で決まり、絶対値の近い固有値の組があると収束は遅い。対称行列に適用すると$A_k$も対称であるから、$A_k$は対角行列に近づく。

\subsection{シフトと減次}
シフトでは、$\lambda_n$の近似値$\sigma$を用いて
\[
    A_k-\sigma I=Q_kR_k,\qquad A_{k+1}=R_kQ_k+\sigma I
\]
と反復する。$R_k=Q_k^{\mathsf T}(A_k-\sigma I)$より$A_{k+1}=Q_k^{\mathsf T}A_kQ_k$であるから、固有値は変わらない。このとき第$n$行の成分の収束率は$|\lambda_n-\sigma|/|\lambda_j-\sigma|$となり、$\sigma$が$\lambda_n$に近ければ0に近いので、第$n$行は非常に速く収束する。$\sigma$には反復中の$A_k$の$(n,n)$成分を用いる。

減次では、第$n$行の非対角成分が十分小さくなった時点で$A_k$の$(n,n)$成分を固有値の1つとして採用し、$A_k$から第$n$行と第$n$列を除いた$(n-1)$次の小行列$B_k$について同じ計算を続ける。第$n$行の非対角成分が0ならば、最後の行で余因子展開して
\[
    \det(zI-A_k)=(z-a_{k;n,n})\det(zI-B_k)
\]
となるので、残りの固有値は$B_k$の固有値で近似できる。行列の次数が下がるため、1反復あたりの計算量も減る。

\subsection{数値解の精度の評価}
数値解の精度は次の量で評価した．
  \begin{itemize}
      \item 残差ノルム

数値解$\tilde{x}$の残差ノルムは$\tilde{x}$をそれが満たす方程式$f(x)=0$があったとき、$\|f(\tilde{x})\|$で定義される。つまり、元の方程式への当てはまり度合を表し、$\tilde{x}$が真の解$x$に等しければ$0$となる。

      \item 相対誤差

      
      真の解 $x$ と数値解 $\tilde{x}$ の差を $x$ の大きさで割った量$\frac{\| x - \tilde{x} \|}{\| x \|}$で定義され、相対的な誤差を示す。
  \end{itemize}

\section{実験方法}

\subsection{実験環境}
計算はPython 3.13.14とNumPy 2.5.3を用いて行い、課題6・7のQR法ではNumba 0.68.0によるJITコンパイルを用いた。図はgnuplotで作成した。計算機の構成は次のとおりである。
\begin{itemize}
    \item OS：Windows 11 Home（バージョン 25H2）
    \item CPU：Intel Core Ultra 5 125U
    \item メモリ：16 GB
\end{itemize}

\subsection{共通の設定}
\begin{itemize}
    \item 行列とベクトルの成分は、NumPyの乱数生成器\texttt{default\_rng}（シード0）による区間$[0,1)$の一様乱数とした。各次数$n$について100個（課題1・2では$A$と$b$の組を100組）生成した。
    \item 課題6・7では、生成した行列$B$を$A=(B+B^{\mathsf T})/2$と対称化して用いた。課題4・5では対称化していない。
    \item 計算時間は\texttt{time.perf\_counter}を用い、数値解（固有値）を求める部分のみを計測した。課題6・7では、Numbaのコンパイル時間を含めないよう、計測の前に1回関数を実行した。
    \item 真の解、固有値、固有ベクトルには、NumPyの\texttt{numpy.linalg.solve}（課題1・2）、\texttt{numpy.linalg.eig}（課題4・5）、\texttt{numpy.linalg.eigh}（課題6・7）で求めた値を用いた。ノルムはすべて2ノルムである。
    \item 課題4・5では、真の固有ベクトルと数値解の固有ベクトルをともに長さ1とし、内積が負のときは真の固有ベクトルの符号を反転して向きをそろえた。課題6・7では、求めた固有値と真の固有値をそれぞれ絶対値の大きい順に並べて対応させた。
\end{itemize}

\subsection{各課題の設定}
\begin{description}
    \item[課題1] 消去法（部分ピボット選択付き）。$n=100,200,400,800$。
    \item[課題2] LU分解による解法（部分ピボット選択付き）。課題1と同じ$A,b$を用いた。
    \item[課題4] べき乗法。$n=50,100,200,400$。初期ベクトル$x_0=(1,\ldots,1)^{\mathsf T}$、収束判定の閾値$\varepsilon=10^{-6}$、反復回数の上限$10^4$。
    \item[課題5] べき乗法を1回行って得た$\mu_0$を$\sigma$とし、$A-\sigma I$を1回だけLU分解したうえで、$y_0/\|y_0\|_2$を初期ベクトルとして逆反復法を行った。収束判定は$\lambda_1$の近似値の相対変化量が$10^{-6}$未満、反復回数の上限は$10^4$。行列は課題4と同じである。
    \item[課題6] QR法（修正Gram--Schmidt法によるQR分解）。$n=10,20,40,80$。収束判定は$A_k$の狭義下三角成分の絶対値の最大値が$10^{-6}$未満、反復回数の上限は$10^4$。
    \item[課題7] シフトと減次を用いたQR法。$\sigma$には各反復で現在の行列の右下成分を用い、最下行の非対角成分の絶対値の最大値が$10^{-6}$未満になった時点で減次した。反復回数の上限は各段で$10^4$。行列は課題6と同じである。
\end{description}

\section{結果}
\subsection{課題１}
 消去法により, 連立 1 次方程式 (2.1) の
解を求める数値実験を行う．n = 100, 200, 400,
800 のそれぞれに対して，n 次の行列 A とベク
トル b のペアをランダムに 100 個生成し，各ペア
に対して連立 1 次方程式 (2.1) の数値解を求めよ．
\begin{equation}
    A x = b
    \tag{2.1}
\end{equation}

\begin{equation}
 \frac{|| x -     \tilde{x}||}{||x||}
 \tag{2.4}
\end{equation}
\begin{equation}
 ||b-A\tilde{x}||
 \tag{2.5}
\end{equation}
\subsubsection{}
各 n に対して，数値解 $\tilde{x}$ の残差ノルム (2.5)，
相対誤差 (2.4)，計算時間をそれぞれプロット
せよ. ただし真の解は不明なので，MATLAB
や Python における数値計算パッケージで得られるものを真の解として用いよ
.

\begin{figure}[htbp]
    \centering
    \panel{work1/residual.png}{a}\hfill
    \panel{work1/relative_error.png}{b}\\[8pt]
    \panel{work1/time.png}{c}
    \caption{消去法（部分ピボット選択付き）で求めた連立 1 次方程式 (2.1) の数値解 $\tilde{x}$ の (a) 残差ノルム (2.5)，(b) 相対誤差 (2.4)，(c) 計算時間 $t$．横軸は次数 $n$ で両対数表示．各 $n=100, 200, 400, 800$ について，成分が区間 $[0,1)$ の一様乱数である $A, b$ を100組生成し，その平均値をプロットした．ノルムは 2 ノルム，真の解 $x$ には NumPy の \texttt{numpy.linalg.solve} の解を用いた．(c) の点線は $n^3$ に比例する参考線である．}
    \label{fig:work1}
\end{figure}

\begin{table}[htbp]
    \centering
    \caption{消去法で求めた数値解 $\tilde{x}$ の残差ノルム (2.5)，相対誤差 (2.4) および計算時間 $t$（各 $n$ で100組の平均値）}
    \label{tab:work1}
    \begin{tabular}{rccc}
        \hline
        $n$ & 残差ノルム & 相対誤差 & $t$/s \\
        \hline
        100 & $8.54\times10^{-14}$ & $3.31\times10^{-14}$ & $1.75\times10^{-3}$ \\
        200 & $4.51\times10^{-13}$ & $8.01\times10^{-14}$ & $5.12\times10^{-3}$ \\
        400 & $1.96\times10^{-12}$ & $2.99\times10^{-13}$ & $7.53\times10^{-2}$ \\
        800 & $6.66\times10^{-12}$ & $7.13\times10^{-13}$ & $5.05\times10^{-1}$ \\
        \hline
    \end{tabular}
\end{table}

\subsubsection{}
数値解の精度と次数 n の関係について考察
せよ．

残差ノルムをみると、すべて$10^{-11}$を下回っている。また、$n=100$と$n=800$の2点から求めた両対数グラフの傾きは
\[
    \frac{\log(6.66\times10^{-12}/8.54\times10^{-14})}{\log(800/100)} \approx 2.10
\]
であり、残差ノルムはおおむね$n^2$程度で増加する傾向がある。

\subsubsection{}
計算時間と次数 n の関係について考察せよ．

計算時間は$n=100$で$1.75\times10^{-3}$\,s、$n=800$で$5.05\times10^{-1}$\,sであり（表\ref{tab:work1}）、2点から求めた両対数グラフの傾きは
\[
    \frac{\log(5.05\times10^{-1}/1.75\times10^{-3})}{\log(800/100)} \approx 2.72
\]
である。前進消去の演算回数は$O(n^3)$であるから傾きは3になると予想されるが、それより小さい。隣り合う$n$の間で求めた傾きは$n=100\to200$で1.55、$n=400\to800$で2.74であり、$n$が大きいほど3に近づいている。$n$が小さいときは、Pythonのループの処理など$n^3$に比例しない部分の時間の割合が大きいためと考えられる。ただし$n=200\to400$の傾きは3.88と3を上回っており、単調な変化にはなっていない。

\clearpage
\subsection{課題２}
課題 1 と同じ内容を，LU 分解による解法に替えて行え．

\subsubsection{}
各 n に対して，数値解 $\tilde{x}$ の残差ノルム (2.5)，
相対誤差 (2.4)，計算時間をそれぞれプロット
せよ. ただし真の解は不明なので，MATLAB
や Python における数値計算パッケージで得られるものを真の解として用いよ．

\begin{figure}[htbp]
    \centering
    \panel{work2/residual.png}{a}\hfill
    \panel{work2/relative_error.png}{b}\\[8pt]
    \panel{work2/time.png}{c}
    \caption{LU 分解（部分ピボット選択付き）による解法で求めた連立 1 次方程式 (2.1) の数値解 $\tilde{x}$ の (a) 残差ノルム (2.5)，(b) 相対誤差 (2.4)，(c) 計算時間 $t$．横軸は次数 $n$ で両対数表示．各 $n=100, 200, 400, 800$ について，成分が区間 $[0,1)$ の一様乱数である $A, b$ を100組生成し，その平均値をプロットした．ノルムは 2 ノルム，真の解 $x$ には NumPy の \texttt{numpy.linalg.solve} の解を用いた．(c) の点線は $n^3$ に比例する参考線である．}
    \label{fig:work2}
\end{figure}

\begin{table}[htbp]
    \centering
    \caption{LU分解による解法で求めた数値解 $\tilde{x}$ の残差ノルム (2.5)，相対誤差 (2.4) および計算時間 $t$（各 $n$ で100組の平均値）}
    \label{tab:work2}
    \begin{tabular}{rccc}
        \hline
        $n$ & 残差ノルム & 相対誤差 & $t$/s \\
        \hline
        100 & $8.17\times10^{-14}$ & $3.28\times10^{-14}$ & $2.42\times10^{-3}$ \\
        200 & $4.46\times10^{-13}$ & $8.06\times10^{-14}$ & $7.51\times10^{-3}$ \\
        400 & $1.86\times10^{-12}$ & $3.00\times10^{-13}$ & $4.77\times10^{-2}$ \\
        800 & $6.81\times10^{-12}$ & $7.16\times10^{-13}$ & $5.11\times10^{-1}$ \\
        \hline
    \end{tabular}
\end{table}

\subsubsection{}
数値解の精度と次数 n の関係について考察せよ．

残差ノルムをみると、すべて$10^{-11}$を下回っている。また、$n=100$と$n=800$の2点から求めた両対数グラフの傾きは
\[
    \frac{\log(6.81\times10^{-12}/8.17\times10^{-14})}{\log(800/100)} \approx 2.13
\]
であり、残差ノルムはおおむね$n^2$程度で増加する傾向がある。
\subsubsection{}
計算時間と次数 n の関係について考察せよ．

計算時間は$n=100$で$2.42\times10^{-3}$\,s、$n=800$で$5.11\times10^{-1}$\,sであり（表\ref{tab:work2}）、2点から求めた両対数グラフの傾きは
\[
    \frac{\log(5.11\times10^{-1}/2.42\times10^{-3})}{\log(800/100)} \approx 2.57
\]
である。LU分解の演算回数は$O(n^3)$、前進代入と後退代入の演算回数は$O(n^2)$であるから、$n$が大きいときの傾きは3になると予想される。隣り合う$n$の間で求めた傾きは$n=100\to200$で1.63、$n=200\to400$で2.67、$n=400\to800$で3.42であり、$n$が大きいほど傾きが大きくなる点は課題1と同じである。

\clearpage
\subsection{課題３}
連立 1 次方程式に対する消去法と LU 分解による解法を，
数値解の精度と計算速度の観点から比較せよ．
その際，理論的な計算量（加減算，乗算，除算それぞれの回数）についても考察せよ．

\begin{table}[htbp]
    \centering
    \footnotesize\setlength{\tabcolsep}{2pt}
    \caption{消去法（課題1）と LU 分解による解法（課題2）の比較．同じ100組の $A, b$ に対する平均値}
    \label{tab:work3}
    \begin{tabular}{rcccccc}
        \hline
        & \multicolumn{2}{c}{残差ノルム (2.5)} & \multicolumn{2}{c}{相対誤差 (2.4)} & \multicolumn{2}{c}{$t$/s} \\
        $n$ & 消去法 & LU分解 & 消去法 & LU分解 & 消去法 & LU分解 \\
        \hline
        100 & $8.54\times10^{-14}$ & $8.17\times10^{-14}$ & $3.31\times10^{-14}$ & $3.28\times10^{-14}$ & $1.75\times10^{-3}$ & $2.42\times10^{-3}$ \\
        200 & $4.51\times10^{-13}$ & $4.46\times10^{-13}$ & $8.01\times10^{-14}$ & $8.06\times10^{-14}$ & $5.12\times10^{-3}$ & $7.51\times10^{-3}$ \\
        400 & $1.96\times10^{-12}$ & $1.86\times10^{-12}$ & $2.99\times10^{-13}$ & $3.00\times10^{-13}$ & $7.53\times10^{-2}$ & $4.77\times10^{-2}$ \\
        800 & $6.66\times10^{-12}$ & $6.81\times10^{-12}$ & $7.13\times10^{-13}$ & $7.16\times10^{-13}$ & $5.05\times10^{-1}$ & $5.11\times10^{-1}$ \\
        \hline
    \end{tabular}
\end{table}

精度については、表\ref{tab:work3}より、残差ノルムの差は5\%以内、相対誤差の差は1\%以内であり、両者はほぼ同じ精度である。LU分解による解法は、消去法の行基本変形を行列$L$、$U$および置換$\pi$として記録したものであり、部分ピボット選択を含めて同じ演算を行うためと考えられる。

理論的な計算量を求める。消去法では、前進消去の第$k$段（$k=1,\ldots,n-1$）で、除算を$n-k$回、乗算と減算をそれぞれ$(n-k)(n-k+1)$回行う。後退代入では、乗算と減算をそれぞれ$n(n-1)/2$回、除算を$n$回行う。したがって
\begin{align*}
    \text{乗算}=\text{減算}&=\sum_{m=1}^{n-1}m(m+1)+\frac{n(n-1)}{2}=\frac{n(n-1)(2n+5)}{6},\\
    \text{除算}&=\frac{n(n-1)}{2}+n=\frac{n(n+1)}{2}
\end{align*}
である。LU分解による解法では、LU分解の第$k$段で除算を$n-k$回、乗算と減算をそれぞれ$(n-k)^2$回行う。前進代入では$L$の対角成分が1なので除算はなく、乗算と減算をそれぞれ$n(n-1)/2$回行う。後退代入は消去法と同じである。したがって
\begin{align*}
    \text{乗算}=\text{減算}&=\sum_{m=1}^{n-1}m^2+n(n-1)=\frac{n(n-1)(2n+5)}{6},\\
    \text{除算}&=\frac{n(n+1)}{2}
\end{align*}
であり、両者の演算回数は一致する。主要項は乗算・減算ともに$n^3/3$である。なお、本実験の実装では第$k$段で左側の0の列も含めて行全体を更新しているため、実際の乗算と減算の回数はいずれも約$n^3/2$回である。

計算速度については、表\ref{tab:work3}より、$n=100,200$では消去法の方が速く（LU分解の0.72倍、0.68倍）、$n=400$ではLU分解の方が速く（消去法がLU分解の1.58倍）、$n=800$ではほぼ同じ（0.99倍）であった。演算回数が同じであるため、この差は行交換の方法や前進代入をPythonのループで行うかどうかといった実装の違いや、計測のばらつきによるものと考えられる。一方、LU分解は同じ$A$に対して異なる$b$を解く場合に分解を1回行えばよいという利点がある。

\clearpage
\subsection{課題４}
べき乗法により，行列 $A$ の第 1 固有値と第 1 固有ベクトルを求める数値実験を行う．
$n = 50, 100, 200, 400$ のそれぞれに対して，$n$ 次行列 $A$ をランダムに 100 個生成し，
各 $A$ の第 1 固有値と第 1 固有ベクトルを数値的に求めよ．
\begin{equation}
 \frac{|\lambda_i - \tilde{\lambda}_i|}{|\lambda_i|}
 \tag{2.10}
\end{equation}
\begin{equation}
 \frac{\| v_i - \tilde{v}_i \|}{\| v_i \|}
 \tag{2.11}
\end{equation}
\begin{equation}
 \| \tilde{\lambda}_i \tilde{v}_i - A \tilde{v}_i \|
 \tag{2.12}
\end{equation}

注意 2.5.1　べき乗法や QR 法の収束性は，与えられた行列の固有値の絶対値がすべて相異なる場合に保証される．
しかし，ランダムな実行列は共役な複素固有値を含む可能性があり，その場合にはアルゴリズムの収束が保証されない．
そこで，生成した実行列 $A$ にその転置 $A^T$ を加えるなどして実対称行列を生成し，
それを用いて実験するなどの工夫をすると良い．

\subsubsection{}
各 $n$ に対して，数値的に求めた第 1 固有値 $\tilde{\lambda}_1$ と第 1 固有ベクトル $\tilde{v}_1$ に対する
固有方程式の残差ノルム (2.12)，$\tilde{\lambda}_1$ の相対誤差 (2.10)，$\tilde{v}_1$ の相対誤差 (2.11)，
計算時間，収束までに要した反復回数をそれぞれプロットせよ．

\begin{figure}[htbp]
    \centering
    \panel{work4/residual.png}{a}\hfill
    \panel{work4/relative_error_lambda.png}{b}\\[8pt]
    \panel{work4/relative_error_x.png}{c}\hfill
    \panel{work4/time.png}{d}\\[8pt]
    \panel{work4/iteration.png}{e}
    \caption{べき乗法で求めた第 1 固有値 $\tilde{\lambda}_1$ と第 1 固有ベクトル $\tilde{v}_1$ の (a) 残差ノルム (2.12)，(b) $\tilde{\lambda}_1$ の相対誤差 (2.10)，(c) $\tilde{v}_1$ の相対誤差 (2.11)，(d) 計算時間 $t$，(e) 反復回数．横軸は次数 $n$ で，(a)--(d) は両対数，(e) は横軸のみ対数で表示．各 $n=50, 100, 200, 400$ について，成分が区間 $[0,1)$ の一様乱数である行列 $A$ を100個生成した．点は各行列に対する値，赤線は各 $n$ での平均値である．初期ベクトルは $x_0=(1,\ldots,1)^{\mathsf T}$，第 $k$ 反復での $\lambda_1$ の近似値を $\mu_k$ として，収束判定は $|\mu_k-\mu_{k-1}|/|\mu_k|<10^{-6}$ とした．真値には NumPy の \texttt{numpy.linalg.eig} で求めた値を用いた．(d) の点線は $n^2$ に比例する参考線である．}
    \label{fig:work4}
\end{figure}

\begin{table}[htbp]
    \centering
    \footnotesize\setlength{\tabcolsep}{3pt}
    \caption{べき乗法で求めた $\tilde{\lambda}_1, \tilde{v}_1$ の残差ノルム (2.12)，相対誤差 (2.10), (2.11)，計算時間 $t$ および反復回数（各 $n$ で100個の平均値）}
    \label{tab:work4}
    \begin{tabular}{rccccc}
        \hline
        & 残差ノルム & $\tilde{\lambda}_1$ の相対誤差 & $\tilde{v}_1$ の相対誤差 & & \\
        $n$ & (2.12) & (2.10) & (2.11) & $t$/s & 反復回数 \\
        \hline
        50 & $2.38\times10^{-6}$ & $9.27\times10^{-8}$ & $1.47\times10^{-8}$ & $2.36\times10^{-4}$ & 6.60 \\
        100 & $3.16\times10^{-6}$ & $6.31\times10^{-8}$ & $3.82\times10^{-9}$ & $8.82\times10^{-5}$ & 6.02 \\
        200 & $9.74\times10^{-6}$ & $9.71\times10^{-8}$ & $4.29\times10^{-9}$ & $1.51\times10^{-4}$ & 5.54 \\
        400 & $5.67\times10^{-6}$ & $2.83\times10^{-8}$ & $1.13\times10^{-9}$ & $3.31\times10^{-4}$ & 5.10 \\
        \hline
    \end{tabular}
\end{table}

\subsubsection{}
$\tilde{\lambda}_1, \tilde{v}_1$ の精度と次数 $n$ の関係について考察せよ．

残差ノルムをみると、すべて$10^{-5}$を下回っている。また、$n=50$と$n=400$の2点から求めた両対数グラフの傾きは
\[
    \frac{\log(5.67\times10^{-6}/2.38\times10^{-6})}{\log(400/50)} \approx 0.42
\]
であり、残差ノルムの$n$に対する増加は緩やかである。ただし$n=200$で最大となっており、単調には増加していない。
\subsubsection{}
計算時間と次数 $n$ の関係について考察せよ．

$n=50$の計算時間の平均値は$2.36\times10^{-4}$\,sであるが、これは図\ref{fig:work4}(d)に見られるように1つだけ$1.44\times10^{-2}$\,sかかった試行が平均を押し上げているためであり、中央値は$8.78\times10^{-5}$\,sである。この試行は全体で最初の計算であり、初回実行時の処理によるものと考えられる。そこで$n=100$と$n=400$の2点から両対数グラフの傾きを求めると
\[
    \frac{\log(3.31\times10^{-4}/8.82\times10^{-5})}{\log(400/100)} \approx 0.95
\]
である。べき乗法の1反復は行列とベクトルの積であり演算回数は$O(n^2)$、反復回数は5から7回程度でほぼ一定なので傾きは2になると予想されるが、それより小さい。本実験の$n$の範囲では、1反復あたりの演算よりもPythonの処理などの$n$によらない時間の割合が大きいためと考えられる。また、反復回数は$n$とともに6.60回から5.10回へ減少している。べき乗法の収束率は$|\lambda_2|/|\lambda_1|$であるから（第2.4.1節）、$n$が大きいほどこの比が小さくなっているためと考えられる。

\clearpage
\subsection{課題５}
べき乗法と逆反復法を組み合わせて，第 1 固有値と第 1 固有ベクトルを数値的に求めるための算法として，
べき乗法より高速なものをつくれ．
つくった算法とべき乗法を精度と速度の面から比較せよ．

\begin{figure}[htbp]
    \centering
    \panel{work5/residual.png}{a}\hfill
    \panel{work5/relative_error_lambda.png}{b}\\[8pt]
    \panel{work5/relative_error_x.png}{c}\hfill
    \panel{work5/time.png}{d}\\[8pt]
    \panel{work5/iteration.png}{e}
    \caption{課題 5 の算法（べき乗法を 1 回行って得た $\mu_0$ を $\sigma$ とし，$A-\sigma I$ を LU 分解して逆反復法を行う）で求めた $\tilde{\lambda}_1, \tilde{v}_1$ の (a) 残差ノルム (2.12)，(b) $\tilde{\lambda}_1$ の相対誤差 (2.10)，(c) $\tilde{v}_1$ の相対誤差 (2.11)，(d) 計算時間 $t$，(e) 反復回数．横軸は次数 $n$ で，(a)--(d) は両対数，(e) は横軸のみ対数で表示．各 $n=50, 100, 200, 400$ について，成分が区間 $[0,1)$ の一様乱数である行列 $A$ を100個生成した（課題 4 と同じ行列）．点は各行列に対する値，赤線は各 $n$ での平均値である．収束判定は $\tilde{\lambda}_1$ の相対変化量が $10^{-6}$ 未満とした．真値には NumPy の \texttt{numpy.linalg.eig} で求めた値を用いた．(d) の点線は $n^3$ に比例する参考線である．}
    \label{fig:work5}
\end{figure}

\begin{table}[htbp]
    \centering
    \footnotesize\setlength{\tabcolsep}{3pt}
    \caption{課題5の算法で求めた $\tilde{\lambda}_1, \tilde{v}_1$ の残差ノルム (2.12)，相対誤差 (2.10), (2.11)，計算時間 $t$ および反復回数（各 $n$ で100個の平均値）}
    \label{tab:work5}
    \begin{tabular}{rccccc}
        \hline
        & 残差ノルム & $\tilde{\lambda}_1$ の相対誤差 & $\tilde{v}_1$ の相対誤差 & & \\
        $n$ & (2.12) & (2.10) & (2.11) & $t$/s & 反復回数 \\
        \hline
        50 & $2.43\times10^{-6}$ & $2.64\times10^{-8}$ & $9.34\times10^{-8}$ & $3.36\times10^{-3}$ & 4.40 \\
        100 & $5.13\times10^{-6}$ & $1.48\times10^{-8}$ & $9.99\times10^{-8}$ & $6.37\times10^{-3}$ & 3.77 \\
        200 & $5.44\times10^{-6}$ & $9.93\times10^{-9}$ & $5.18\times10^{-8}$ & $1.06\times10^{-2}$ & 3.32 \\
        400 & $6.97\times10^{-6}$ & $5.65\times10^{-9}$ & $3.31\times10^{-8}$ & $6.94\times10^{-2}$ & 3.07 \\
        \hline
    \end{tabular}
\end{table}

\begin{table}[htbp]
    \centering
    \footnotesize\setlength{\tabcolsep}{3pt}
    \caption{べき乗法（課題4）と課題5の算法の精度の比較．同じ100個の $A$ に対する平均値}
    \label{tab:work5_acc}
    \begin{tabular}{rcccccc}
        \hline
        & \multicolumn{2}{c}{残差ノルム (2.12)} & \multicolumn{2}{c}{$\tilde{\lambda}_1$ の相対誤差 (2.10)} & \multicolumn{2}{c}{$\tilde{v}_1$ の相対誤差 (2.11)} \\
        $n$ & べき乗法 & 課題5 & べき乗法 & 課題5 & べき乗法 & 課題5 \\
        \hline
        50 & $2.38\times10^{-6}$ & $2.43\times10^{-6}$ & $9.27\times10^{-8}$ & $2.64\times10^{-8}$ & $1.47\times10^{-8}$ & $9.34\times10^{-8}$ \\
        100 & $3.16\times10^{-6}$ & $5.13\times10^{-6}$ & $6.31\times10^{-8}$ & $1.48\times10^{-8}$ & $3.82\times10^{-9}$ & $9.99\times10^{-8}$ \\
        200 & $9.74\times10^{-6}$ & $5.44\times10^{-6}$ & $9.71\times10^{-8}$ & $9.93\times10^{-9}$ & $4.29\times10^{-9}$ & $5.18\times10^{-8}$ \\
        400 & $5.67\times10^{-6}$ & $6.97\times10^{-6}$ & $2.83\times10^{-8}$ & $5.65\times10^{-9}$ & $1.13\times10^{-9}$ & $3.31\times10^{-8}$ \\
        \hline
    \end{tabular}
\end{table}

\begin{table}[htbp]
    \centering
    \footnotesize\setlength{\tabcolsep}{3pt}
    \caption{べき乗法（課題4）と課題5の算法の計算時間 $t$ と反復回数の比較．同じ100個の $A$ に対する平均値}
    \label{tab:work5_spd}
    \begin{tabular}{rcccc}
        \hline
        & \multicolumn{2}{c}{$t$/s} & \multicolumn{2}{c}{反復回数} \\
        $n$ & べき乗法 & 課題5 & べき乗法 & 課題5 \\
        \hline
        50 & $2.36\times10^{-4}$ & $3.36\times10^{-3}$ & 6.60 & 4.40 \\
        100 & $8.82\times10^{-5}$ & $6.37\times10^{-3}$ & 6.02 & 3.77 \\
        200 & $1.51\times10^{-4}$ & $1.06\times10^{-2}$ & 5.54 & 3.32 \\
        400 & $3.31\times10^{-4}$ & $6.94\times10^{-2}$ & 5.10 & 3.07 \\
        \hline
    \end{tabular}
\end{table}

精度については、表\ref{tab:work5_acc}より、課題5の算法の$\tilde{\lambda}_1$の相対誤差はべき乗法の約$1/3.5$から$1/10$と小さい一方、$\tilde{v}_1$の相対誤差はべき乗法の約6倍から29倍と大きい。残差ノルムはどちらも$10^{-6}$程度で同程度である。相対誤差はいずれも収束判定の閾値$10^{-6}$より小さい。

速度については、表\ref{tab:work5_spd}より、反復回数はべき乗法の6.60回から5.10回に対して4.40回から3.07回と少ないが、計算時間はべき乗法の約14倍から210倍となった。課題5の算法では、最初に$A-\sigma I$のLU分解（演算回数$O(n^3)$）を行い、さらに各反復で前進代入と後退代入をPythonのループで行っている。一方、べき乗法は各反復で行列とベクトルの積を1回計算するだけである。べき乗法がもともと5から7回程度で収束するため、反復回数を減らしてもLU分解と代入の時間を補えなかったと考えられる。したがって本実験では、課題の要求であるべき乗法より高速な算法にはなっていない。

\clearpage
\subsection{課題６}
QR 法により，行列 $A$ の全ての固有値を求める数値実験を行う．
$n = 10, 20, 40, 80$ のそれぞれに対して，$n$ 次行列 $A$ をランダムに 100 個生成し，
各 $A$ の固有値 $\lambda_1, \ldots, \lambda_n$ を数値的に求めよ．

\subsubsection{}
各 $n$ に対して，数値的に求めた固有値 $\tilde{\lambda}_1, \ldots, \tilde{\lambda}_n$ に対する
固有方程式の残差ノルムの最大値
\[
 \max_{1 \leq i \leq n} \| (A - \tilde{\lambda}_i I) v_i \|,
\]
相対誤差の最大値
\[
 \max_{1 \leq i \leq n} \frac{|\lambda_i - \tilde{\lambda}_i|}{|\lambda_i|},
\]
計算時間，収束までに要した反復回数をそれぞれプロットせよ．
ただし真の固有値 $\lambda_i$ と固有ベクトル $v_i$ は不明なので，
MATLAB や Python における数値計算パッケージで得られるものを真の固有値および固有ベクトルとして用いよ．

\begin{figure}[htbp]
    \centering
    \panel{work6/residual.png}{a}\hfill
    \panel{work6/relative_error.png}{b}\\[8pt]
    \panel{work6/time.png}{c}\hfill
    \panel{work6/iteration.png}{d}\\[8pt]
    \caption{QR 法（修正 Gram--Schmidt 法による QR 分解，シフト・減次なし）で求めた固有値 $\tilde{\lambda}_1,\ldots,\tilde{\lambda}_n$ の (a) 残差ノルムの最大値 $\max_i\|(A-\tilde{\lambda}_iI)v_i\|_2$，(b) 相対誤差の最大値 $\max_i|\lambda_i-\tilde{\lambda}_i|/|\lambda_i|$，(c) 計算時間 $t$，(d) 反復回数．横軸は次数 $n$ で両対数表示．各 $n=10, 20, 40, 80$ について，成分が区間 $[0,1)$ の一様乱数である行列 $B$ を100個生成し，$A=(B+B^{\mathsf T})/2$ と対称化して用いた．真の固有値 $\lambda_i$ と固有ベクトル $v_i$ には NumPy の \texttt{numpy.linalg.eigh} で求めた値を用いた．点は各行列に対する値，赤線は各 $n$ での平均値，(c) の点線は $n^3$ に比例する参考線である．収束判定は第 $k$ 反復の行列 $A_k$ の狭義下三角成分の絶対値の最大値が $10^{-6}$ 未満，反復回数の上限は $10^4$ とした．}
    \label{fig:work6}
\end{figure}

\begin{table}[htbp]
    \centering
    \footnotesize\setlength{\tabcolsep}{3pt}
    \caption{QR法で求めた固有値に対する残差ノルムの最大値，相対誤差の最大値，計算時間 $t$ および反復回数（各 $n$ で100個の平均値）}
    \label{tab:work6}
    \begin{tabular}{rcccc}
        \hline
        $n$ & 残差ノルムの最大値 & 相対誤差の最大値 & $t$/s & 反復回数 \\
        \hline
        10 & $6.03\times10^{-7}$ & $7.08\times10^{-7}$ & $5.03\times10^{-2}$ & $1.40\times10^{3}$ \\
        20 & $3.68\times10^{-2}$ & $4.07\times10^{-2}$ & $4.84\times10^{-1}$ & $2.64\times10^{3}$ \\
        40 & $2.17\times10^{-1}$ & $1.94\times10^{-1}$ & $2.33\times10^{1}$ & $6.56\times10^{3}$ \\
        80 & $1.36$ & $6.15\times10^{-1}$ & $8.39\times10^{2}$ & $9.45\times10^{3}$ \\
        \hline
    \end{tabular}
\end{table}

\subsubsection{}
$\tilde{\lambda}_1, \ldots, \tilde{\lambda}_n$ の精度と次数 $n$ の関係について考察せよ．

残差ノルムの最大値をみると、$n=10$では$6.03\times10^{-7}$であるが$n=80$では$1.36$に達している。$n=10$と$n=80$の2点から求めた両対数グラフの傾きは
\[
    \frac{\log(1.36/6.03\times10^{-7})}{\log(80/10)} \approx 7.03
\]
であり、残差ノルムの最大値は$n$とともに急激に増加している。
\subsubsection{}
計算時間と次数 $n$ の関係について考察せよ．

計算時間の平均値は$n=10$で$5.03\times10^{-2}$\,s、$n=80$で$8.39\times10^{2}$\,sであり（表\ref{tab:work6}）、2点から求めた両対数グラフの傾きは
\[
    \frac{\log(8.39\times10^{2}/5.03\times10^{-2})}{\log(80/10)} \approx 4.67
\]
と3を大きく上回る。QR法の1反復ではQR分解と行列積を行い、演算回数は$O(n^3)$である。これに加えて反復回数も$n$とともに$1.40\times10^3$回から$9.45\times10^3$回へ増えるため、計算時間は$n^3$より速く増加すると考えられる。なお、$n=80$では図\ref{fig:work6}(c)に$10^4$\,sを超える点が1つあり、平均値を押し上げている。中央値で同様に傾きを求めると約4.2である。

\clearpage
\subsection{課題７}
シフトと減次を用いて QR 法を高速化せよ．
つくった算法と元々の QR 法を精度と速度の面から比較せよ．

\begin{figure}[htbp]
    \centering
    \panel{work7/residual.png}{a}\hfill
    \panel{work7/relative_error.png}{b}\\[8pt]
    \panel{work7/time.png}{c}\hfill
    \panel{work7/iteration.png}{d}\\[8pt]
    \caption{シフトと減次を用いた QR 法で求めた固有値 $\tilde{\lambda}_1,\ldots,\tilde{\lambda}_n$ の (a) 残差ノルムの最大値 $\max_i\|(A-\tilde{\lambda}_iI)v_i\|_2$，(b) 相対誤差の最大値 $\max_i|\lambda_i-\tilde{\lambda}_i|/|\lambda_i|$，(c) 計算時間 $t$，(d) 反復回数．横軸は次数 $n$ で両対数表示．各 $n=10, 20, 40, 80$ について，成分が区間 $[0,1)$ の一様乱数である行列 $B$ を100個生成し，$A=(B+B^{\mathsf T})/2$ と対称化して用いた（課題 6 と同じ行列）．真の固有値 $\lambda_i$ と固有ベクトル $v_i$ には NumPy の \texttt{numpy.linalg.eigh} で求めた値を用いた．点は各行列に対する値，赤線は各 $n$ での平均値，(c) の点線は $n^3$ に比例する参考線である．シフト量 $\sigma$ には各反復で現在の行列の右下成分を用い，最下行の非対角成分の絶対値の最大値が $10^{-6}$ 未満になった時点で減次した．反復回数は全段の QR 反復の合計である．}
    \label{fig:work7}
\end{figure}

\begin{table}[htbp]
    \centering
    \footnotesize\setlength{\tabcolsep}{3pt}
    \caption{シフトと減次を用いたQR法で求めた固有値に対する残差ノルムの最大値，相対誤差の最大値，計算時間 $t$ および反復回数（各 $n$ で100個の平均値）}
    \label{tab:work7}
    \begin{tabular}{rcccc}
        \hline
        $n$ & 残差ノルムの最大値 & 相対誤差の最大値 & $t$/s & 反復回数 \\
        \hline
        10 & $2.40\times10^{-12}$ & $2.35\times10^{-11}$ & $1.53\times10^{-4}$ & 18.1 \\
        20 & $3.82\times10^{-12}$ & $2.42\times10^{-11}$ & $8.87\times10^{-4}$ & 38.1 \\
        40 & $1.07\times10^{-11}$ & $7.90\times10^{-11}$ & $5.60\times10^{-3}$ & 74.7 \\
        80 & $2.31\times10^{-11}$ & $1.36\times10^{-10}$ & $4.71\times10^{-2}$ & 148 \\
        \hline
    \end{tabular}
\end{table}

\begin{table}[htbp]
    \centering
    \footnotesize\setlength{\tabcolsep}{3pt}
    \caption{QR法（課題6）とシフトと減次を用いたQR法（課題7）の精度の比較．同じ100個の $A$ に対する平均値}
    \label{tab:work7_acc}
    \begin{tabular}{rcccc}
        \hline
        & \multicolumn{2}{c}{残差ノルムの最大値} & \multicolumn{2}{c}{相対誤差の最大値} \\
        $n$ & 課題6 & 課題7 & 課題6 & 課題7 \\
        \hline
        10 & $6.03\times10^{-7}$ & $2.40\times10^{-12}$ & $7.08\times10^{-7}$ & $2.35\times10^{-11}$ \\
        20 & $3.68\times10^{-2}$ & $3.82\times10^{-12}$ & $4.07\times10^{-2}$ & $2.42\times10^{-11}$ \\
        40 & $2.17\times10^{-1}$ & $1.07\times10^{-11}$ & $1.94\times10^{-1}$ & $7.90\times10^{-11}$ \\
        80 & $1.36$ & $2.31\times10^{-11}$ & $6.15\times10^{-1}$ & $1.36\times10^{-10}$ \\
        \hline
    \end{tabular}
\end{table}

\begin{table}[htbp]
    \centering
    \footnotesize\setlength{\tabcolsep}{3pt}
    \caption{QR法（課題6）とシフトと減次を用いたQR法（課題7）の計算時間 $t$ と反復回数の比較．同じ100個の $A$ に対する平均値}
    \label{tab:work7_spd}
    \begin{tabular}{rcccc}
        \hline
        & \multicolumn{2}{c}{$t$/s} & \multicolumn{2}{c}{反復回数} \\
        $n$ & 課題6 & 課題7 & 課題6 & 課題7 \\
        \hline
        10 & $5.03\times10^{-2}$ & $1.53\times10^{-4}$ & $1.40\times10^{3}$ & 18.1 \\
        20 & $4.84\times10^{-1}$ & $8.87\times10^{-4}$ & $2.64\times10^{3}$ & 38.1 \\
        40 & $2.33\times10^{1}$ & $5.60\times10^{-3}$ & $6.56\times10^{3}$ & 74.7 \\
        80 & $8.39\times10^{2}$ & $4.71\times10^{-2}$ & $9.45\times10^{3}$ & 148 \\
        \hline
    \end{tabular}
\end{table}

精度については、表\ref{tab:work7_acc}より、シフトと減次を用いたQR法では、すべての$n$で残差ノルムの最大値の平均値が$2.31\times10^{-11}$以下、相対誤差の最大値の平均値が$1.36\times10^{-10}$以下であり、課題6のQR法より大幅に改善した。図\ref{fig:work7}(d)より、反復回数が上限$10^4$に達した行列はなく、すべての行列で収束している。

速度については、表\ref{tab:work7_spd}より、計算時間は課題6のQR法の約330分の1（$n=10$）から約$1.8\times10^4$分の1（$n=80$）となった。反復回数は$18.1$回から$148$回であり、$n$で割るとどの$n$でも約1.8回から1.9回となる。つまり固有値1つあたり2回程度のQR反復で第$m$行が収束している。シフトによって第$m$行の収束率がほぼ0になり（第2.4.5節）、さらに減次によって行列の次数が小さくなっていくためと考えられる。

\clearpage
\section{考察・議論}

本実験では、連立1次方程式の直接法2種類と、行列固有値問題の反復法4種類を実装し、精度と計算時間を調べた。

連立1次方程式については、消去法とLU分解による解法の演算回数は一致し、精度もほぼ同じであった（課題3）。いずれも相対誤差は$10^{-12}$未満であり、$n=800$までの範囲で十分な精度が得られた。計算時間の理論的な次数は$O(n^3)$であるが、本実験の範囲ではPythonの処理などの影響により、両対数グラフの傾きは3より小さかった。

固有値問題については、べき乗法は課題4の行列に対して5回から7回程度の反復で収束した。課題4・5では注意2.5.1の対称化を行っていないが、成分がすべて正の行列では、絶対値が最大の固有値は正の実数であり、他の固有値より絶対値が真に大きい（Perronの定理）。そのため本実験の行列ではべき乗法の収束の条件が満たされていたと考えられる。課題5の算法は反復回数を減らしたものの、LU分解とPythonのループによる代入に時間がかかり、べき乗法より遅くなった。代入をNumPyなどの三角行列用の関数で行うことや、$\sigma$を反復の途中で$\lambda_1$により近い値に取り替えることが改善策として考えられるが、後者では$A-\sigma I$のLU分解をやり直す必要がある。

QR法については、シフトと減次を用いない場合、絶対値の近い固有値の組をもつ行列で収束が遅く、$n=80$では大半の行列で反復回数が上限に達した。シフトと減次を用いることで、すべての行列が固有値1つあたり2回程度の反復で収束し、精度と計算時間の両方が大幅に改善した。
\section{結論}

消去法とLU分解による解法は、演算回数・精度ともに同等であり、$n=800$まで相対誤差$10^{-12}$未満の解が得られた。べき乗法は本実験の行列に対して少ない反復回数で収束したが、べき乗法と逆反復法を組み合わせた算法は、LU分解と代入の時間のためべき乗法より遅くなった。QR法はシフトと減次を用いない場合には$n$が大きいと収束しない行列が多かったが、シフトと減次を用いることで、すべての行列で固有値を$10^{-10}$程度の相対誤差で求められ、計算時間も大幅に短縮された。
\section{参考文献}

\begin{thebibliography}{9}
\bibitem{text} 川越大輔, 数値線形代数, 数理工学実験 (2026), 第2章.
\end{thebibliography}
\end{document}
